\subsubsection*{4.2.1.}
先考虑特例 $d=2$ 且 $X$ 为 $\mathbb P^r$ 中的线性子簇. 只需在 $\check{\mathbb P}^{N}$ 中找到一条不与 $F$ 相交的直线(参见 3.2.4). 选取 $\mathbb P^r$ 上的齐次坐标 $X_0,\ldots,X_r$, 使 $X$ 由
\[
X_{n+1}=\cdots=X_r=0,
\]
定义, 并在 $\mathbb P^1$ 上选取齐次坐标 $(\lambda,\mu)$.

当 $n$ 为偶数, $n=2k$, 取二次超曲面束
\begin{equation}\tag{4.2.2}
H_{(\lambda,\mu)}
=\lambda X_0^2+(\lambda X_k+\mu X_0)X_{2k}
 +\sum_{i=1}^{k-1}(\lambda X_i+\mu X_{i+1})X_{k+i}.
\end{equation}
当 $n$ 为奇数, $n=2k+1$, 取束
\begin{equation}\tag{4.2.3}
H_{(\lambda,\mu)}
=(\lambda X_1+\mu X_0)X_0
 +\sum_{i=1}^{k}(\lambda X_{2i}+\mu X_i)X_{2i+1}.
\end{equation}
在两种情形中, 均可检验当 $(\lambda,\mu)\ne(0,1)$ 时 $X\cap H_{(\lambda,\mu)}$ 光滑, 而 $X\cap H_{(0,1)}$ 恰有一个奇点, 且为普通二次奇点. 因此此束不与 $F$ 相交.

\subsubsection*{4.2.4.}
设 $\operatorname{Btg}^{d}(X)$ 为
\[
\check{\mathbb P}^{N}\times
\bigl((X\times X)-\Delta(X\times X)\bigr)
\]
中由三元组 $(H,x_1,x_2)$ 构成的子簇, 其中 $H$ 为 $\mathbb P^r$ 中在 $x_1$ 和 $x_2$ 处均与 $X$ 相切的 $d$ 次超曲面, 且 $x_1\ne x_2$. $\operatorname{Btg}^{d}(X)$ 在 $\operatorname{pr}_1$ 下的像就是命题 4.2 中的集合 $F_2$, 故只需证明
\begin{equation}\tag{4.2.5}
\dim\operatorname{Btg}^{d}(X)\leq N-2
\quad\text{若 $d\geq3$, 或 $d=2$ 且 $X$ 非线性.}
\end{equation}

为此, 设 $\operatorname{Lin}(X\times X)$ 为 $(X\times X)-\Delta(X\times X)$ 中由
\begin{equation}\tag{4.2.6}
\begin{aligned}
\operatorname{Lin}(X\times X)
&=\bigl\{(x_1,x_2)\in X\times X\mid x_1\ne x_2,\\[-1ex]
&\hspace{8em}\text{由 $x_1,x_2$ 张成的直线在 $x_1$ 处与 $X$ 相切}
\bigr\}\\
&=\bigl\{(x_1,x_2)\in X\times X\mid x_1\ne x_2,\\[-1ex]
&\hspace{8em}x_2\in X\cap\operatorname{Tg}_X(x_1)\bigr\}.
\end{aligned}
\end{equation}
给出的子簇. 显然:

\medskip
\noindent(4.2.7) \underline{若 $X$ 不是 $\mathbb P^r$ 的线性子簇, 则}
\[
\dim\operatorname{Lin}(X\times X)<2\dim X.
\]
因此, 为证明 (4.2.5), 只需证明态射
\begin{equation}\tag{4.2.8}
\operatorname{Btg}^{d}(X)
\xrightarrow{\ \operatorname{pr}_{2,3}\ }
(X\times X)-\Delta(X\times X),
\end{equation}
的纤维视为 $\check{\mathbb P}^{N}$ 的子簇时满足
\begin{equation}\tag{4.2.9}
\left\{
\begin{array}{ll}
\text{余维至少为 $2\dim X+1$}
 &\text{在 $\operatorname{Lin}(X\times X)$},\\[1ex]
\text{余维至少为 $2\dim X+2$}
 &\text{在 $(X\times X)-\Delta(X\times X)
                 -\operatorname{Lin}(X\times X)$.}
\end{array}
\right.
\end{equation}
取 $L_i$ 为 $X$ 在 $x_i$ 处的切空间, $i=1,2$, 则陈述 (4.2.9) 由下一命题得出.